\documentclass{../pset}

% User Information
\course{APS106H1: Fundamentals of Computer Programming}
\assignment{Problem Set 2}


\begin{document}

\begin{problem}
Write a function that takes in a three-digit integer and returns the "flipped" version (that is, with its digits reversed). For example:
 
\begin{lstlisting}[language={}]
$ Enter a positive number less than 1000: 177
That number flipped is 771
\end{lstlisting}
Your function should take the user-input integer as an argument and return the flipped integer. That means all your input and output with the user should be in your "main" code, not in the function.

\textit{Hints:}
\begin{itemize}
\item  Consider using the arithmetic operators \lstinline|//| and \lstinline|%| to solve this problem.
\item Make sure it works if the user enters a one- or two-digit integer. For example, 37 flipped is 730, and 8 flipped is 800.
\end{itemize}
\end{problem}



\begin{problem}

Write a function that takes in a positive integer and prints the minimum number of quarters, dimes, nickels, and pennies needed to make up that amount.  
For example:
\begin{lstlisting}[language={}]
$ Enter a number of cents: 67
2 quarter(s), 1 dime(s), 1 nickel(s), 2 penny(ies)
\end{lstlisting}


The prompt for the user should be outside the function but the printing should be \textbf{inside} the function.
In case you are not familiar with Canadian currency: 
\begin{itemize}
\item  a quarter is 25 cents
\item  a dime is 10 cents
\item  a nickel is 5 cents
\item a penny is 1 cent
\end{itemize}
Use the minimum number of coins that you can by first maximizing the number of quarters, then the number of dimes, etc

\end{problem}


\begin{problem}
 Write a function that takes floating-point number and then prints two values: the number truncated to the first decimal place and the number rounded to the first decimal place. For example:

\begin{lstlisting}[language={}]
$ Enter a number: 4.158
Truncated to one decimal place: 4.1
Rounded to one decimal place: 4.2
\end{lstlisting}
The prompt for the user should be \textbf{outside} the function but the printing should be inside the function

\textit{Hint:} For truncating, there are multiple ways to achieve this. Think about using \lstinline|int|, or alternatively, you can also think about using the \href{https://docs.python.org/3/library/math.html}{\lstinline|math|} module.
\end{problem}


\begin{problem}

Write a  trigonometry function. Ask a user for an angle, specified in \textbf{degrees}. Then print the sine, cosine, and tangent of that angle. Round to 3 decimal places. For example:
\begin{lstlisting}[language={}]
$ Enter an angle in degrees:  60
sin(60.0) is 0.866
cos(60.0) is 0.5
tan(60.0) is 1.732
\end{lstlisting}



Printing should occur \textbf{inside} the function. You can use \href{https://docs.python.org/3/library/math.html}{\lstinline|math|} module, where applicable.
\end{problem}

\begin{problem}


Write a function that takes a positive integer input less than 100,000 and returns an integer corresponding to the number of digits in the number. Use this function in a program that prompts the user for an integer and prints the number of digits in the user's input. All input and printing to the user should be \textbf{outside} the function. For example:
\begin{lstlisting}[language={}]
$ Please input a number less than 100,000:  1234
You number has 4 digits.
\end{lstlisting}



\textbf{Note,}  try to solve this problem through multiple approaches:
\begin {itemize}
\item Think about how to use string \lstinline|str| and \lstinline|len| function to approach it.
\item Also think about solving it mathematically by transforming the integer. Hint, the transformation function in question can be found in the \href{https://docs.python.org/3/library/math.html}{\lstinline|math|} module.
\item Don't use an if-statement (that is for next week). 
\end{itemize}
\end{problem}

\end{document}